\section{Introduction}


% 1. Context and the Problem
% 2. State of the Art and the Gap
% 3. Proposed approach and contributions
% 4. Paper Structure


%Industry analysts also increasingly recognize sustainability as a core software quality attribute: Gartner predicts that by 2027, around 30\% of large global enterprises will include software sustainability as a formal non-functional requirement in their development processes, up from less than 10 % in 2024

Source code is the fundamental component of software systems, defining the behavior of systems ranging from embedded IoT devices to global communication networks \cite{DiCosmo2017}. While hardware efficiency has improved significantly, the software running on it remains a massive consumer of energy, often written without regard for its environmental footprint due to software bloat and abstraction layers \cite{Wirth1995}. The aggregate energy consumption of inefficient code across billions of devices contributes heavily to global carbon emissions \cite{Belkhir2018, Freitag2021}. To mitigate this, we must move beyond traditional post-execution profiling and develop predictive models capable of estimating the energy consumption of code blocks \textbf{during the development phase}. By quantifying the energy cost of specific constructs before compilation, developers can treat energy as a tangible resource and make informed decisions in real-time, rather than relying solely on retrospective measurements \cite{Pang2016, Pinto2017}.  

\begin{comment}

Software systems are a fundamental component of modern infrastructure, powering everything from global communication networks to the rapidly expanding Internet of Things (IoT). However, this scale comes with a substantial environmental cost. While individual applications may seem negligible, their aggregate energy consumption is massive. Data centers alone account for a significant fraction of global electrical usage \cite{Marantos2023TSUSC, doi:10.1126/science.aba3758}, and the cumulative energy consumption of software on billions of consumer devices contributes heavily to carbon emissions. Consequently, energy efficiency is no longer just a requirement for high-performance computing; it is a necessity for all general-purpose software. Reducing energy consumption helps lower the environmental carbon footprint and significantly reduces operational costs for companies running software at scale\cite{10.1145/2597073.2597097}. 

\end{comment}

Despite this need, developing energy-efficient software remains structurally difficult. Current assessment techniques rely predominantly on \textit{runtime measurements} using hardware power meters or processor-level counters like Intel RAPL \cite{10.1145/2425248.2425252}. This creates a reactive \textit{"measure-after-build"} workflow, where a program must be fully implemented and executed before its energy behavior can be observed. This late-stage feedback is costly; discovering an energy defect after deployment often requires expensive refactoring cycles. Furthermore, existing profiling tools often provide insufficient granularity for effective optimization \cite{7429297}. They typically report energy usage at the process or application level, failing to attribute computation to specific code constructs such as loops or conditional blocks that developers actively manage \cite{Khan2018RAPLinAction, Shivadharshan2024CPPJoules}. This disconnect leaves developers unable to pinpoint the exact source of energy inefficiencies, reducing optimization efforts to trial and error \cite{Pinto2017}. 

To address these limitations, proactive approaches are needed that shift energy assessment from a runtime concern to a design-time quality attribute. Developers currently rely on static analysis tools and linters such as \textit{SonarQube} \cite{Lenarduzzi2020JSS}, \textit{PMD} \cite{AlOmar2023SEET_PMD},\textit{CheckStyle}\footnote{\texttt{Checkstyle}: \url{https://checkstyle.sourceforge.io/}}
, and \textit{SpotBugs}\footnote{\texttt{SpotBugs}: \url{https://spotbugs.github.io/}} to identify bugs and enforce coding standards before code is even compiled. A similar proactive approach is required for energy efficiency. However, there are insufficient tools that can estimate the energy implications of source code structure during the development phase. To the best of our knowledge, no existing approach provides fine-grained, block-level energy estimation purely from source code at design time without requiring code execution or specialized hardware profiling.  This gap forces energy efficiency to remain a retroactive afterthought rather than a proactive design parameter. 

In this paper, we propose \textbf{\textit{EnCoDe}}, a novel methodology for \textbf{\textit{"design-time, fine-grained estimation of energy consumption"}}. Our methodology bridges the gap between static code structure and dynamic energy behavior. To achieve this, we first address the measurement gap by developing \textbf{\textit{PowerLens}}, a custom measurement methodology capable of reliably profiling code blocks at sub-millisecond granularity. This infrastructure allows us to construct a ground-truth dataset of code blocks annotated with precise energy consumption values. Leveraging this dataset, we extract rich structural features from \textbf{\textit{Abstract Syntax Trees (ASTs)}} and train machine learning models to predict block-level energy consumption statically. We explicitly employ classical, low-overhead machine learning approaches (e.g., Random Forests, Gradient Boosting) rather than energy-intensive deep learning models. This alignment with the principles of Data-Centric Green AI~\cite{Verdecchia2022DataCentricGreenAI} ensures that our solution does not ironically consume excessive energy in the pursuit of energy efficiency, maintaining a low carbon footprint for the tool itself. 

The specific contributions of this paper are as follows:
\begin{enumerate}
    \item \textbf{\textit{PowerLens Measurement:}} We present a novel measurement methodology that utilizes execution amplification and temporal synchronization to \textbf{achieve reliable sub-millisecond energy readings}. This allows us to accurately profile microsecond-scale code blocks that are otherwise invisible to standard hardware counters.
    \item \textbf{\textit{Fine-Grained Energy Dataset:}} We conducted an extensive empirical study on executable code blocks extracted from over \textbf{18,000 Python programs}. This study uncovers the linear and non-linear relationships between static code features such as structural complexity, operator density, and energy consumption.
    \item \textbf{\textit{Predictive Modeling and Validation:}} We demonstrate that classical machine learning models trained on these static features can accurately estimate energy consumption at development time. Our results show that EnCoDe achieves a coefficient of determination \textbf{\(R^2\) of 0.75} for regression and an \textbf{accuracy of 80.6\%} for classifying energy hotspots, enabling developers to identify and address inefficiencies early in the software development lifecycle. 
\end{enumerate}

The remainder of this paper is structured as follows. \textbf{Section ~\ref{sec:related}} reviews the related work in software energy measurement and static analysis. \textbf{Section ~\ref{sec:methodology}} detailed the EnCoDe methodology, including the PowerLens methodology and the feature extraction process. \textbf{Section ~\ref{sec:exp-setup}} describes the experimental setup and dataset construction. \textbf{Section ~\ref{sec:research-questions}} formulates the research questions guiding this study. \textbf{Section ~\ref{sec:results}} presents our experimental results and analysis. \textbf{Section ~\ref{sec:Discussion}} discusses the implications of our finding.\textbf{section ~\ref{sec:threats}} addresses the threats to validity. Finally, \textbf{Section ~\ref{sec:conclusion}} documents our conclusions and outlines future work. 


\begin{comment}
    
Software has become an integral part of modern life, running on devices from smartphones and laptops to massive cloud infrastructures. With this ubiquity comes an enormous energy footprint. Data centres alone account for a significant fraction of global electricity usage \cite{doi:10.1126/science.aba3758}, while the cumulative energy consumption of everyday software on personal devices is substantial. As society moves toward sustainability, energy-efficient software has become a pressing need: it not only reduces environmental impact but also lowers operational costs and improves performance for battery-constrained devices \cite{10.1145/2597073.2597097}.

While developers may intend to write energy-efficient code, doing so in practice remains extremely challenging. Current techniques for assessing software energy efficiency rely almost entirely on \emph{runtime measurements}, using hardware power meters or processor-level energy counters such as Intel RAPL \cite{10.1145/2425248.2425252,RAPL}. As a consequence, a program must first be fully implemented and executed on target hardware before its energy behavior can be observed. Once inefficiencies are discovered, developers must refactor the code, a process that is time-consuming, costly, and therefore often neglected in practice. As a result, energy optimization is rarely attempted outside performance-critical domains.

Moreover, the granularity of feedback provided by existing energy profiling tools is often insufficient for effective optimization \cite{7429297}. Runtime profilers typically report aggregate energy consumption at the process level, and even advanced approaches rarely go beyond method-level attribution. This leaves developers with a frustrating question: \emph{which part of my code is responsible for the energy waste?} Without fine-grained insights into energy hotspots, optimization efforts become largely guesswork. If developers cannot prioritize which code regions to improve, the potential for meaningful energy savings is lost.

What is needed is the equivalent of linters and code quality tools that provide actionable feedback to developers at design time, before execution. Just as linters guide developers toward maintainable, secure, and bug-free code by enforcing quality attributes, an analogous tool should guide developers toward energy efficiency by estimating the cost of different code structures. Recent surveys in green software engineering emphasize the need to integrate energy considerations earlier in the software development lifecycle, rather than treating energy as a purely runtime concern \cite{LEE2024111944}. Fine-grained energy measurement combined with predictive models can transform energy efficiency from a late-stage afterthought into an early-stage design concern.

In this paper, we take a step toward that vision. We present a novel framework for \textbf{design-time, fine-grained estimation of energy consumption}. At the core of our approach is a custom-built measurement infrastructure that reliably profiles code blocks at sub-millisecond granularity. This enables us to construct a dataset of code blocks annotated with fine-grained ground-truth energy consumption. We then extract rich structural features from Abstract Syntax Trees (ASTs) and train machine learning models to predict block-level energy consumption statically, from source code alone. Once trained, these models can analyze unseen code without execution, providing actionable feedback about energy-intensive blocks before the program is ever run.

By making energy efficiency measurable and predictable at design time, our work elevates energy awareness to the same level as traditional code quality concerns. Just as linters guide developers in writing secure, readable, and reliable code, EnCoDe guides them in writing energy-efficient software from the outset.
\end{comment}
